%% raw_hw.tex -- Homework 1-3: prompts (verbatim), my graded submissions (HW1, HW2), my own solutions (HW3)
%% Deduped against Part I (d1, d2, d2b, d4, d5 "HW answers at a glance"); proofs removed. Compressed (c5).
\section{Homework 1--3: prompts, my submissions, solutions}
\ixh{homework}%
{\color{figgray}\itshape [Omitted (in Part I: d5 ``HW answers at a glance'', d2/d2b pipelines, d4 templates): HW1 setup, Q1--Q6, most plot code; HW2 setup, Q5, Q7, Q10--Q12, most recipe/workflow/metric code + tables; HW3 setup, Q7, Q11--Q12, model specs/workflows, AUC tables, confusion matrix + metrics.]}

{\color{figgray}\itshape Key: \textbf{Prompt.} = assignment (condensed); \textbf{\color{dblue}My answer.} = my graded submission (HW1, HW2); \textbf{\color{keyc}Answer.} = correction/note (HW1, HW2) or my own solution, run in R 4.5.3 (HW3, unsubmitted).}

%% =====================================================================
\subsection{Homework 1 -- ML main ideas; EDA on mpg}
\ixch{mpg}\ixh{exploratory data analysis (EDA)}%
\begin{qbox}{HW1 Ex1}
\textbf{Prompt.} \ix{histogram} of \ixc{hwy} (highway mpg); describe.
\mine
\fig{binwidth 2, x = Highway MPG, y = Count. Counts (centers 12, 14, \dots, 44): 5, 12, 38, 23, 13, 14, 28, 46, 29, 11, 6, 3, 3, 0, 1, 0, 2.}
Two peaks, $\approx$16--18 and 24--28 mpg; median 24, mean $\approx$23.4. Most cars 15--30 mpg; a few outliers\ixh{outlier} $>40$.
\end{qbox}

\begin{qbox}{HW1 Ex2}
\textbf{Prompt.} \ix{scatterplot}, \ixc{hwy} on x, \ixc{cty} on y. Relationship? Meaning?
\mine
\fig{tight straight upward band from (12, 9) to (44, 35); integer grid, so only 78 distinct points for 234 cars.}
City MPG consistently $<$ highway; fitted slope $\approx$0.7 (+1 hwy mpg $\to$ +0.7 cty mpg). Fewer than 234 dots: rounding to whole numbers $\to$ overlap.\ixh{correlation}
\end{qbox}

\begin{qbox}{HW1 Ex3}
\textbf{Prompt.} Horizontal \ix{bar plot} of \ixc{manufacturer}, ordered by height. Most / fewest cars?
\mine
\begin{rcode}
mpg %>% count(manufacturer, sort = TRUE)  (*@\ixch{\%>\%}\ixch{count()}@*)
\end{rcode}
\begin{routput}
dodge 37, toyota 34, volkswagen 27, ford 25, chevrolet 19,
audi 18, hyundai 14, subaru 14, nissan 13, honda 9, jeep 8,
pontiac 5, land rover 4, mercury 4, lincoln 3   # 15 rows
\end{routput}
\end{qbox}

\begin{qbox}{HW1 Ex4}
\textbf{Prompt.} \ix{boxplot} of \ixc{hwy} by \ixc{cyl}, add points with \ixc{geom\_jitter()} + \ixc{alpha}. Relationship?
\mine
\begin{rcode}
ggplot(mpg, aes(x = factor(cyl), y = hwy)) +  (*@\ixch{factor()}@*)
  geom_boxplot(outlier.shape = NA) +  (*@\ixch{geom\_boxplot()}\ixch{outlier.shape}@*)
  geom_jitter(alpha = 0.3, width = 0.2) +  (*@\ixch{geom\_jitter()}\ixch{alpha}@*)
  labs(x = "Number of Cylinders", y = "Highway MPG")
\end{rcode}
\fig{4 cyl: box 26--31, points 20--44 (three $>40$); 5 cyl: 4 cars, box 28.75--29; 6 cyl: box 19--26, points 17--29; 8 cyl: box 16--19, points 12--26.}
More cylinders $\to$ less chance of full utilization $\to$ lower efficiency. 4cyl widest spread, has the $>40$ outliers. 5cyl too few cars for a reliable box. Groups overlap (some 4cyl worse than some 8cyl), but on average cyl negatively correlated with mpg.
\ans (note) 4-cyl widest \emph{range} (20--44), but widest \emph{box} (IQR) is 6-cyl (7 vs.\ 5). \texttt{outlier.shape = NA} stops outliers being drawn twice (box + jitter).
\end{qbox}

\begin{qbox}{HW1 Ex5}
\textbf{Prompt.} Lower-triangle \ix{correlation matrix} of \texttt{mpg} with \ixc{corrplot}. Which pairs $\pm$ correlated? Sensible / surprising?
\mine
\fig{numeric vars, blue = $+$, red = $-$. year--displ 0.15; cyl--displ 0.93, cyl--year 0.12; cty--displ $-0.80$, cty--year $-0.04$, cty--cyl $-0.81$; hwy--displ $-0.77$, hwy--year 0.00, hwy--cyl $-0.76$, hwy--cty 0.96.}
\texttt{displ}--\texttt{cyl} strongly $+$ (more cylinders = more volume displaced per cycle). Both mpg vars strongly $-$ with both engine-size vars (more volume/mass to heat $\to$ higher fuel burn). Year uncorrelated with mpg --- surprising (EPA tightening), but data only 1999 and 2008.\ixh{Pearson correlation}
\end{qbox}

{\color{figgray}\itshape [Ex6--8: 231 only, from \texttt{homework-1-1.html}; not in my submission. Answers = my recreations, checked in R.]}

\begin{qbox}{HW1 Ex6 (231 only)}
\textbf{Prompt.} Recreate with \ixc{ggplot()} (hint: \ixc{ggthemes}):
\fig{horizontal boxplots, x = Highway MPG, y = Vehicle Class (2seater \dots\ suv), unfilled boxes, black outlier dots, gray jittered points, Google-Docs-style theme. Pickup, suv lowest ($\approx$17); compact/midsize/subcompact highest (medians $\approx$26--27); subcompact widest.}
\ans
\begin{rcode}
ggplot(mpg, aes(x = hwy, y = class)) +
  geom_boxplot() + geom_jitter(alpha = 0.3) +
  theme_gdocs() +  (*@\ixch{theme\_gdocs()}@*)
  labs(x = "Highway MPG", y = "Vehicle Class")
\end{rcode}
\end{qbox}

\begin{qbox}{HW1 Ex7 (231 only)}
\textbf{Prompt.} Recreate:
\fig{default theme; boxplots of \texttt{hwy} by \texttt{class}, fill = \texttt{drv} (4 red, f green, r blue); medians: 2seater r $\approx$25; compact 4 $\approx$25, f $\approx$29 (outliers 37, 44); midsize 4 $\approx$24, f $\approx$27; minivan f $\approx$23; pickup 4 $\approx$17; subcompact 4 26 (flat), f $\approx$29 (to 44), r $\approx$23; suv 4 $\approx$18 (many outliers), r $\approx$17.}
\ans
\begin{rcode}
ggplot(mpg, aes(x = class, y = hwy, fill = drv)) + geom_boxplot()  (*@\ixch{fill}@*)
\end{rcode}
\end{qbox}

\begin{qbox}{HW1 Ex8 (231 only)}
\textbf{Prompt.} Recreate:
\fig{\texttt{hwy} vs.\ \texttt{displ}, points colored by drv; one blue loess curve per drv, no ribbon, linetype by drv, one merged legend: 4 solid, $\approx$27$\to$16 then flat; f short dash, $\approx$33$\to$25 then flat; r long dash, $\approx$25, dip to $\approx$18 at displ 5.3, back to $\approx$25.}
\ans
\begin{rcode}
ggplot(mpg, aes(x = displ, y = hwy)) +
  geom_point(aes(color = drv)) +  (*@\ixch{color}@*)
  geom_smooth(aes(linetype = drv), se = FALSE)  (*@\ixch{geom\_smooth()}\ixch{linetype}\ixch{se = FALSE}@*)
\end{rcode}
\end{qbox}

%% =====================================================================
\subsection{Homework 2 -- linear regression and KNN (abalone)}
\ixch{abalone}\ixh{linear regression}\ixh{k-nearest neighbors (KNN)}%
\textbf{Prompt.} UCI abalone data: 4{,}177 abalone (Tasmania). Age (= \textbf{rings + 1.5}) normally needs cutting the shell and counting rings under a microscope; goal: predict age from easier measurements. Codebook:
{\footnotesize
\begin{tabular}{@{}l l l l@{}}
\toprule
Sex & nominal & -- & M, F, I (infant)\\
Length & cont. & mm & longest shell measurement\\
Diameter & cont. & mm & perpendicular to length\\
Height & cont. & mm & with meat in shell\\
Whole weight & cont. & g & whole abalone\\
Shucked weight & cont. & g & weight of meat\\
Viscera weight & cont. & g & gut weight (after bleeding)\\
Shell weight & cont. & g & after being dried\\
Rings & integer & -- & +1.5 = age in years\\
\bottomrule
\end{tabular}}
{\color{figgray}\itshape [CSV names: \ixc{type}, \ixc{longest\_shell}, \ixc{diameter}, \ixc{height}, \ixc{whole\_weight}, \ixc{shucked\_weight}, \ixc{viscera\_weight}, \ixc{shell\_weight}, \ixc{rings}.]}

\begin{qbox}{HW2 Q1}
\textbf{Prompt.} Add \ixc{age} (= rings + 1.5); describe its distribution.
\mine
\begin{rcode}
head(abalone)  (*@\ixch{head()}@*)
\end{rcode}
\begin{lstlisting}[style=cs/rout,basicstyle=\ttfamily\scriptsize]
  type  longest_shell diameter height whole_weight shucked_weight viscera_weight
1 M             0.455    0.365  0.095        0.514         0.224          0.101
2 M             0.35     0.265  0.09         0.226         0.0995         0.0485
3 F             0.53     0.42   0.135        0.677         0.256          0.142
4 M             0.44     0.365  0.125        0.516         0.216          0.114
5 I             0.33     0.255  0.08         0.205         0.0895         0.0395
6 I             0.425    0.3    0.095        0.352         0.141          0.0775
# i 3 more variables: shell_weight <dbl>, rings <dbl>, age <dbl>
\end{lstlisting}
\begin{rcode}
abalone %>% summarise(n = n(), mean = mean(age), median = median(age),  (*@\ixch{summarise()}\ixch{n()}@*)
  sd = sd(age), min = min(age), q1 = quantile(age, 0.25),  (*@\ixch{quantile()}@*)
  q3 = quantile(age, 0.75), max = max(age))
\end{rcode}
\begin{routput}
      n  mean median    sd   min    q1    q3   max
1  4177  11.4   10.5  3.22   2.5   9.5  12.5  30.5
\end{routput}
\fig{histogram (binwidth 1): unimodal, peak $\approx$690 at 10.5 (11.5 $\approx$630, 9.5 $\approx$570, 12.5 $\approx$490), long thin right tail to 30.5 (bars barely visible $>24$). Boxplot: box 9.5--12.5, median 10.5, whiskers 5.5--16.5, outliers 2.5--4.5 and 17.5--30.5.}
Max $\approx3\times$ median, but IQR only 3 years $\to$ tightly bound around median.
\ans (check) Correct: \texttt{boxplot.stats()} flags 278/4177 = 6.66\% as outliers (outside $9.5-1.5\cdot3=5$, $12.5+1.5\cdot3=17$).
\end{qbox}

\begin{qbox}{HW2 Q2}
\textbf{Prompt.} Train/test split with \ix{stratified sampling}; choose proportions.
\mine 85/15 (Andrew Ng's split, I think): 3548 train rows, 629 test rows --- both adequate.\ixh{training/testing split}\ixh{data splitting}
\end{qbox}

\begin{qbox}{HW2 Q3}
\textbf{Prompt.} Recipe on \textbf{training} data: \texttt{age} \textasciitilde{} all predictors \textbf{except} \ixc{rings}. \emph{Why not use \texttt{rings}?}
\mine
\begin{rcode}
abalone_recipe %>% prep() %>% bake(new_data = NULL) %>% head()  (*@\ixch{prep()}\ixch{bake()}\ixch{new\_data}@*)
\end{rcode}
\begin{lstlisting}[style=cs/rout,basicstyle=\ttfamily\scriptsize]
  longest_shell diameter height whole_weight shucked_weight viscera_weight
1        -1.62    -1.54  -1.41        -1.27          -1.22          -1.29
2        -0.828   -1.09  -1.05        -0.975         -0.985         -0.943
3        -1.41    -1.29  -1.29        -1.10          -1.19          -1.29
4        -0.496   -0.537 -0.817       -0.715         -0.599         -0.519
5        -0.621   -0.537 -0.817       -0.627         -0.554         -0.587
6        -2.66    -2.60  -1.99        -1.61          -1.50          -1.51
# i 8 more variables: shell_weight <dbl>, age <dbl>, type_I <dbl>,
#   type_M <dbl>, type_I_x_shucked_weight <dbl>, type_M_x_shucked_weight <dbl>,
#   longest_shell_x_diameter <dbl>, shucked_weight_x_shell_weight <dbl>
\end{lstlisting}
\ans (note) \ix{target leakage}: \texttt{age} $=$ \texttt{rings} $+1.5$ exactly $\to$ perfect fit ($R^2=1$), uninformative; and \texttt{rings} is the hard-to-obtain measurement we want to avoid.
\end{qbox}

\begin{qbox}{HW2 Q4}
\textbf{Prompt.} Linear regression object, \texttt{"lm"} engine.
\mine \texttt{Linear Regression Model Specification (regression) / Computational engine: lm}
\end{qbox}

\begin{qbox}{HW2 Q6}
\textbf{Prompt.} For LM and \ix{KNN}: two separate workflows (empty workflow $\to$ add model $\to$ add Q3 recipe); fit both to training.
\mine (\texttt{lm\_fit \%>\% extract\_fit\_parsnip() \%>\% tidy()}, then \texttt{knn\_fit}):
\begin{routput}
   term                       estimate std.error statistic  p.value
 1 (Intercept)                  11.5      0.0367   312.    0
 2 longest_shell                 0.466    0.281      1.66  9.67e- 2
 3 diameter                      2.20     0.308      7.16  1.00e-12
 4 height                        0.262    0.0689     3.81  1.43e- 4
 5 whole_weight                  4.98     0.390     12.8   1.71e-36
 6 shucked_weight               -4.26     0.247    -17.3   3.07e-64
 7 viscera_weight               -1.02     0.155     -6.54  7.22e-11
 8 shell_weight                  1.61     0.209      7.72  1.52e-14
 9 type_I                       -0.919    0.111     -8.26  1.99e-16
10 type_M                       -0.223    0.101     -2.21  2.71e- 2
11 type_I_x_shucked_weight       0.468    0.0843     5.54  3.16e- 8
12 type_M_x_shucked_weight       0.237    0.106      2.23  2.59e- 2
13 longest_shell_x_diameter     -2.79     0.394     -7.06  1.94e-12
14 shucked_weight_x_shell_weight -0.149   0.201     -0.742 4.58e- 1
\end{routput}
\ixh{coefficient}\ixh{p-value}\ixh{standard error}\ixh{t-statistic}%
\begin{routput}
kknn::train.kknn(formula = ..y ~ ., data = data,
  ks = min_rows(7, data, 5))
Type of response variable: continuous
minimal mean absolute error: 1.654218
Minimal mean squared error: 5.598762
Best kernel: optimal
Best k: 7
\end{routput}
\end{qbox}

\begin{qbox}{HW2 Q8}
\textbf{Prompt.} With \ixc{yardstick}, for LM and KNN: metric set of $R^2$, RMSE, MAE;\ixh{R-squared}\ixh{RMSE}\ixh{MAE} \ixc{augment()} on \textbf{testing} data (predictions + observed age); apply metrics, report, interpret $R^2$.
\mine (\texttt{head()} of \texttt{augment(...) \%>\% select(age, .pred)}, LM and KNN side by side):
\begin{routput}
    age .pred(LM) .pred(KNN)
1   8.5  9.51      9.31
2  11.5 12.5      12.0
3  13.5 12.9      13.4
4   8.5 10.5      10.9
5  12.5  9.91      8.68
6  13.5 12.0      11.7
\end{routput}
\fig{Predicted vs.\ observed test age (4--23), facets KNN (k = 7) and LM, red dashed $y=x$: both follow the line for young abalone, fall below it (under-predict) above age $\approx$15; KNN more spread.}
\end{qbox}

\begin{qbox}{HW2 Q9}
\textbf{Prompt.} Which model did better on testing? Why? Surprised?\ixh{model comparison}
\mine
\begin{rcode}
bind_rows(
  abalone_metrics(lm_test_results, truth = age, estimate = .pred) %>%
    mutate(model = "Linear Regression", data = "Testing"),
  abalone_metrics(knn_test_results, truth = age, estimate = .pred) %>%
    mutate(model = "KNN (k = 7)", data = "Testing"),
  augment(lm_fit, new_data = abalone_train) %>%
    abalone_metrics(truth = age, estimate = .pred) %>%
    mutate(model = "Linear Regression", data = "Training"),
  augment(knn_fit, new_data = abalone_train) %>% ...  # same, "KNN (k = 7)"
) %>%
  select(model, data, .metric, .estimate) %>%  (*@\ixch{.estimate}@*)
  pivot_wider(names_from = .metric, values_from = .estimate) %>%  (*@\ixch{pivot\_wider()}@*)
  arrange(data, model)  (*@\ixch{arrange()}@*)
\end{rcode}
\begin{routput}
  model             data       rsq  rmse   mae
1 KNN (k = 7)       Testing  0.483  2.21  1.62
2 Linear Regression Testing  0.585  1.95  1.44
3 KNN (k = 7)       Training 0.811  1.46  1.02
4 Linear Regression Training 0.552  2.18  1.57
\end{routput}
\end{qbox}

%% =====================================================================
\subsection{Homework 3 -- binary classification (Titanic)}
\ixch{titanic}\ixh{classification}\ixh{logistic regression}\ixh{linear discriminant analysis (LDA)}\ixh{quadratic discriminant analysis (QDA)}%
\textbf{Prompt.} Part of the Kaggle Titanic data; predict which passengers survived. Codebook:\par\noindent
{\footnotesize
\begin{tabular}{@{}l p{0.78\linewidth}@{}}
survival / sex / Age: & Survival / Sex / Age in years\\
pclass: & Ticket class: 1 = 1st, 2 = 2nd, 3 = 3rd\\
sibsp / parch: & \# siblings+spouses / parents+children aboard\\
ticket / fare / cabin: & Ticket number / Passenger fare / Cabin number\\
embarked: & Port: C = Cherbourg, Q = Queenstown, S = Southampton\\
\end{tabular}}
{\color{figgray}\itshape [CSV: \ixc{passenger\_id}, \texttt{survived} (Yes/No), \texttt{pclass}, \texttt{name}, \ixc{sex}, \ixc{age}, \ixc{sib\_sp}, \ixc{parch}, \texttt{ticket}, \ixc{fare}, \ixc{cabin}, \ixc{embarked}; 891 rows. \texttt{set.seed(3435)}.]}

\begin{qbox}{HW3 Q1}
\textbf{Prompt.} Split, stratified on \texttt{survived}; verify row counts; note issues (e.g.\ missing data) in training. Why stratify here?\ixh{stratified sampling}\ixh{training/testing split}
\begin{rcode}
colSums(is.na(titanic_train))  (*@\ixch{is.na()}\ixch{colSums()}@*)
\end{rcode}
\ans \ix{missing data}: \texttt{age} 128/623 (20.5\%; \ixc{step\_impute\_linear()} in Q4), \texttt{cabin} 481 (77\%, too sparse), \texttt{embarked} 2. \texttt{fare} right-skewed (max 512).
\end{qbox}

\begin{qbox}{HW3 Q2}
\textbf{Prompt.} Describe distribution of \texttt{survived} (training). Percent stacked bar charts of \texttt{survived} (x) with \texttt{fill = sex}, then \texttt{fill = pclass}: good predictors? Why percent stacked vs.\ traditional stacked?\ixh{percent stacked bar chart}\ixh{bar plot}
\begin{rcode}
titanic_train %>% count(survived) %>% mutate(prop = n / sum(n))  (*@\ixch{count()}\ixch{mutate()}@*)
titanic_test %>% count(survived) %>% mutate(prop = n / sum(n))
\end{rcode}
\begin{routput}
  survived  n  prop       survived  n  prop   # train | test
1 Yes     239 0.384     1 Yes     103 0.384
2 No      384 0.616     2 No      165 0.616
\end{routput}
Chart numbers (\% within each bar) + survival rate per group (same code with \texttt{pclass}):
\begin{rcode}
titanic_train %>% count(survived, sex) %>%
  group_by(survived) %>% mutate(prop = n / sum(n))  (*@\ixch{group\_by()}@*)
titanic_train %>% group_by(sex) %>%
  summarise(n = n(), surv_rate = mean(survived == "Yes"))  (*@\ixch{summarise()}@*)
\end{rcode}
\begin{routput}
  survived sex        n  prop       sex        n surv_rate
1 Yes      female   156 0.653     1 female   214     0.729
2 Yes      male      83 0.347     2 male     409     0.203
3 No       female    58 0.151
4 No       male     326 0.849
  survived pclass     n  prop       pclass     n surv_rate
1 Yes      1         97 0.406     1 1        154     0.630
2 Yes      2         61 0.255     2 2        128     0.477
3 Yes      3         81 0.339     3 3        341     0.238
4 No       1         57 0.148
5 No       2         67 0.174
6 No       3        260 0.677
\end{routput}
\end{qbox}

\begin{qbox}{HW3 Q3}
\textbf{Prompt.} \ix{correlation matrix} of all continuous variables (training); visualize; which predictors correlated, which direction?
\begin{rcode}
cor_titanic <- titanic_train %>%
  select(age, sib_sp, parch, fare) %>%  (*@\ixch{select()}@*)
  cor(use = "pairwise.complete.obs")  (*@\ixch{cor()}\ixch{use = "pairwise.complete.obs"}@*)
round(cor_titanic, 3)
corrplot(cor_titanic, method = "number", type = "lower")  (*@\ixch{corrplot()}@*)
\end{rcode}
\begin{routput}
          age sib_sp  parch  fare
age     1.000 -0.327 -0.173 0.113
sib_sp -0.327  1.000  0.433 0.145
parch  -0.173  0.433  1.000 0.254
fare    0.113  0.145  0.254 1.000
\end{routput}
\ans (\texttt{passenger\_id} = ID, \texttt{pclass} = factor, so excluded.) \ix{Pearson correlation}: \texttt{sib\_sp}--\texttt{parch} $+$ (0.43, families travel together); \texttt{age}--\texttt{sib\_sp} $-$ ($-0.33$, children travel with siblings); \texttt{age}--\texttt{parch} weak $-$; \texttt{fare} weak $+$ with \texttt{parch}, \texttt{sib\_sp} (fare is per ticket, families share one) and \texttt{age}.
\end{qbox}

\begin{qbox}{HW3 Q4}
\textbf{Prompt.} Recipe (training) for \texttt{survived} from pclass, sex, age, sib\_sp, parch, fare; impute \texttt{age} with \ixc{step\_impute\_linear()}; \textbf{dummy} encode categoricals with \ixc{step\_dummy()}; interactions sex:fare and age:fare (find the step in tidymodels docs).\ixh{imputation}\ixh{linear imputation}\ixh{dummy variable}\ixh{interaction term}
\begin{rcode}
titanic_recipe <- recipe(survived ~ pclass + sex + age + sib_sp +  (*@\ixch{recipe()}@*)
                           parch + fare, data = titanic_train) %>%
  step_impute_linear(age,  (*@\ixch{step\_impute\_linear()}@*)
                     impute_with = imp_vars(pclass, sex, sib_sp,  (*@\ixch{impute\_with}\ixch{imp\_vars()}@*)
                                            parch, fare)) %>%
  step_dummy(all_nominal_predictors()) %>%  (*@\ixch{step\_dummy()}\ixch{all\_nominal\_predictors()}@*)
  step_interact(terms = ~ starts_with("sex"):fare + age:fare)  (*@\ixch{step\_interact()}\ixch{starts\_with()}@*)
baked <- titanic_recipe %>% prep() %>% bake(new_data = NULL)  (*@\ixch{prep()}\ixch{bake()}@*)
baked; sum(is.na(baked$age))
\end{rcode}
\begin{routput}
# A tibble: 623 x 10
     age sib_sp parch   fare survived pclass_X2 pclass_X3 sex_male
 1    22      1     0   7.25 No               0         1        1
 2    35      0     0   8.05 No               0         1        1
 3    54      0     0  51.9  No               0         0        1
 4     2      3     1  21.1  No               0         1        1
 5    39      1     5  31.3  No               0         1        1
 6    14      0     0   7.85 No               0         1        0
# ... (4 more rows shown in R)
# i 2 more variables: sex_male_x_fare <dbl>, fare_x_age <dbl>
[1] 0
\end{routput}
\ans Imputation = linear model of \texttt{age} on the other 5 predictors (training rows with age observed). \ixc{step\_dummy()}: \texttt{sex} $\to$ \ixc{sex\_male} (female = \ix{reference level}); \texttt{pclass} $\to$ \ixc{pclass\_X2}, \ixc{pclass\_X3} (class 1 = reference). 1 outcome, 6 inputs $\to$ 9 predictors, no missing age.
\end{qbox}

\begin{qbox}{HW3 Q5}
\textbf{Prompt.} \textbf{\ix{logistic regression}}, \texttt{"glm"} engine; workflow + recipe; \texttt{fit()} on training.
\ans (\texttt{log\_fit \%>\% extract\_fit\_parsnip() \%>\% tidy()}):
\begin{routput}
   term             estimate std.error statistic  p.value
 1 (Intercept)     -4.40      0.683       -6.45  1.13e-10
 2 age              0.0629    0.0136       4.62  3.77e- 6
 3 sib_sp           0.437     0.132        3.30  9.52e- 4
 4 parch            0.151     0.153        0.989 3.23e- 1
 5 fare            -0.00116   0.0107      -0.108 9.14e- 1
 6 pclass_X2        1.25      0.363        3.46  5.48e- 4
 7 pclass_X3        2.44      0.382        6.39  1.62e-10
 8 sex_male         2.15      0.303        7.09  1.32e-12
 9 sex_male_x_fare  0.0139    0.00836      1.66  9.65e- 2
10 fare_x_age      -0.000360  0.000206    -1.75  8.03e- 2
\end{routput}
``Yes'' is the first level $\to$ \texttt{glm} models $P(\text{``No''})$: \ix{log odds} of \emph{dying}. \texttt{sex\_male} $+2.15$: male odds of dying $e^{2.15}\approx8.6\times$ female's (at \texttt{fare} = 0, others fixed); \texttt{pclass\_X3} $+2.44$: $e^{2.44}\approx11.5\times$ vs.\ 1st class; \texttt{age} $+0.063$/yr. Significant (\ix{p-value} $<0.05$): age, sib\_sp, pclass, sex; interactions borderline ($p\approx0.08$--$0.10$).\ixh{coefficient}\ixh{odds ratio}
\end{qbox}

\begin{qbox}{HW3 Q6}
\textbf{Prompt.} Repeat Q5 with \ix{linear discriminant analysis (LDA)}, \texttt{"MASS"} engine.
\ans (\texttt{lda\_fit}; group means re-wrapped):
\begin{routput}
lda(..y ~ ., data = data)
Prior probabilities of groups:
      Yes        No
0.3836276 0.6163724
Group means:
         age    sib_sp     parch     fare pclass_X2 pclass_X3
Yes 28.74378 0.4393305 0.4267782 48.81424 0.2552301 0.3389121
No  29.64519 0.5807292 0.3463542 22.99299 0.1744792 0.6770833
     sex_male sex_male_x_fare fare_x_age
Yes 0.3472803        12.91229   1568.232
No  0.8489583        19.63018    703.509
Coefficients of linear discriminants (LD1):
age 0.0358259965     sib_sp 0.2231961671   parch 0.0876726369
fare 0.0032525962    pclass_X2 0.8724616082
pclass_X3 1.6983328837   sex_male 1.9274805785
sex_male_x_fare 0.0020072010   fare_x_age -0.0001487374
\end{routput}
Survivors paid more (48.8 vs.\ 23.0), less often male (0.35 vs.\ 0.85), less often 3rd class (0.34 vs.\ 0.68). LD1: positive weights on sex\_male, pclass\_X3 point toward ``No''.\ixh{discriminant function}\ixh{prior probability}
\end{qbox}

\begin{qbox}{HW3 Q8}
\textbf{Prompt.} Repeat Q5 with \emph{k}-nearest neighbors\ixh{k-nearest neighbors (KNN)}, \texttt{"kknn"} engine; choose \emph{k}.
\ans $k = 10$; \texttt{knn\_fit \%>\% extract\_fit\_parsnip()}:
\begin{routput}
Type of response variable: nominal
Minimal misclassification: 0.1990369
Best kernel: optimal
Best k: 10
\end{routput}
($\approx$80\% leave-one-out training accuracy.) \ix{KNN} is distance-based, needs comparable scales; recipe has no \ixc{step\_normalize()}, but \texttt{kknn} standardizes internally (\texttt{scale = TRUE} default), so \texttt{fare} does not dominate the \ix{Euclidean distance}.\ixh{standardization}
\end{qbox}

\begin{qbox}{HW3 Q9}
\textbf{Prompt.} With \ixc{predict()} + \ixc{bind\_cols()}, predict all 4 models on \textbf{training} data; assess by \textbf{area under the ROC curve}.\ixh{AUC}\ixh{ROC curve}
\ans \texttt{head(log\_train)} (\texttt{predict(..., type = "prob")} + \texttt{bind\_cols(survived)}):
\begin{routput}
  .pred_Yes .pred_No survived
1    0.114     0.886 No
2    0.0835    0.916 No
3    0.310     0.690 No
4    0.115     0.885 No
5    0.0219    0.978 No
6    0.755     0.245 No
\end{routput}
Extra: training accuracy (default \texttt{type = "class"} $\to$ \ixc{.pred\_class}, \ixc{accuracy()}): Logistic 0.814, LDA 0.793, QDA 0.778, KNN 0.891.
\end{qbox}

\begin{qbox}{HW3 Q10}
\textbf{Prompt.} Report each model's \textbf{testing} AUC; which is highest? For the best model: \ix{confusion matrix} (+ visualize) and ROC curve plot. Compare its training vs.\ testing AUC; why different?\ixh{model evaluation}\ixh{testing data}
\begin{rcode}
test_pred <- function(f) predict(f, titanic_test, type = "prob") %>%
  bind_cols(predict(f, titanic_test)) %>%
  bind_cols(titanic_test %>% select(survived))
log_test <- test_pred(log_fit); lda_test <- test_pred(lda_fit)
qda_test <- test_pred(qda_fit); knn_test <- test_pred(knn_fit)
\end{rcode}
Extra: testing accuracy Logistic 0.832, LDA 0.806, QDA 0.810, KNN 0.843.
\fig{LDA test ROC (x = 1 $-$ specificity, y = sensitivity): sens $\approx$0.70 at 0.05, $\approx$0.85 at 0.18, $\approx$0.90 from 0.25 to 0.6, then (1, 1); AUC 0.881.}
\end{qbox}

\textbf{Prompt (231 only).} Binary problem: $p$ = P(class 1), $1-p$ = P(class 0). The \emph{\ix{logistic function}} (``inverse logit'') is the CDF of the logistic distribution, mapping real \emph{z} to $(0, 1)$. {\color{figgray}\itshape [Q11--Q12 + answers: Part I.]}
